\chapter{Vision Transformer}
\label{cap:vit}

Il Vision Transformer \parencite{dosovitskiy2020image} applica l'architettura del
\cref{cap:transformer} alle immagini, trattandole come sequenze di regioni
rettangolari. Il capitolo si concentra sul \emph{patch embedding}, che è il
componente modificato negli esperimenti 02 e 03.

\paragraph{Notazione.} L'immagine ha forma $\shape{\dbatch, \nch, H, W}$; le patch
sono quadrate di lato $p$ e non sovrapposte, dunque
$\npatch = (H/p)(W/p)$ e la patch appiattita ha dimensione
$\dpatch = p^2\nch$. La dimensione dei token è $\dvit$
(\code{embed\_dim}). Con i valori del repository: $H = W = 28$, $p = 7$,
$\nch = 3$, quindi $\npatch = 16$ e $\dpatch = 147$; $\dvit = 10$
nell'esperimento~03 e $\dvit = 8$ nel~02.

\section{Dalle patch ai token}

L'attenzione opera su sequenze, mentre un'immagine è una griglia di pixel. Il
Vision Transformer partiziona l'immagine in regioni disgiunte e tratta ciascuna
come un token, così che ogni pixel contribuisca a esattamente un token.
L'estrazione usa due applicazioni di \code{unfold}, seguite da una \code{permute}
che sposta l'asse dei canali in ultima posizione prima dell'appiattimento. Dopo
la \code{view} finale, le tre componenti RGB dello stesso pixel sono contigue nel
vettore di $\dpatch$ elementi. Il repository usa immagini quadrate e verifica che
la dimensione del lato sia divisibile per $p$.

\begin{definizione}[Patch embedding]
La proiezione delle patch è la mappa $\R^{\dpatch} \to \R^{\dvit}$ che trasforma
ogni patch appiattita in un token del Transformer. Nella formulazione base è una
trasformazione affine, $\text{PatchEmbed}(x) = xW + b$ con
$W \in \R^{\dpatch \times \dvit}$.
\end{definizione}

Con $\dpatch = 147$ e $\dvit = 10$ la proiezione ha $147 \times 10 + 10 =
\num{1480}$ parametri, verificati nel campo \code{patch\_embed} di
\file{params.json} della variante \code{vanilla}, e il rapporto di compressione è
$147/10 = 14{,}7$.

Prima dell'encoder, il \emph{patch embedding} riduce ogni vettore da 147 a 8
componenti nell'esperimento 02 e a 10 nel 03. Il VQC riceve quindi una
rappresentazione prodotta da uno strato classico addestrabile, non i pixel
grezzi. Le tre varianti dell'ablazione sono descritte nella
\cref{sec:implementazioni-vit} e il loro design nel \cref{cap:esperimenti}.

\section{Token di classificazione ed embedding posizionale}
\label{sec:cls}

Un parametro appreso $c \in \R^{\shape{1,1,\dvit}}$, il token \code{[CLS]}, è
concatenato in testa alla sequenza e non è associato a una patch. Attraverso
l'attenzione, la sua rappresentazione in uscita dall'encoder raccoglie
informazioni dalla sequenza ed è usata per la classificazione. Un'alternativa,
non valutata in questi esperimenti, è il \emph{mean pooling} sulle
rappresentazioni finali delle patch.

L'attenzione è equivariante rispetto alle permutazioni della sequenza, quindi
senza informazione posizionale la griglia di patch sarebbe per il modello un
insieme non ordinato. L'embedding posizionale è un parametro appreso di forma
$\shape{1, \npatch + 1, \dvit}$, dove il $+1$ tiene conto del token \code{[CLS]};
viene sommato ai token e seguito da dropout. Diversamente dal caso testuale,
l'embedding è appreso in una sola dimensione pur codificando posizioni
bidimensionali: la corrispondenza fra indice di sequenza e coordinata nella
griglia resta implicita nell'ordine prodotto dall'estrazione delle patch.

I due esperimenti inizializzano questi parametri su scale diverse: deviazione
standard $0{,}02$ nel 03 e unitaria nel 02.

\begin{limite}
La diversa scala di inizializzazione, insieme alle dimensioni dei modelli,
esclude confronti diretti fra i risultati numerici dei due esperimenti, ma non
incide sui rispettivi confronti appaiati.
\end{limite}

\section{Encoder e testa di classificazione}

L'encoder usa l'implementazione standard di PyTorch con due strati e due teste,
configurazione pre-norm e attivazione GELU. La dimensione della rete feed-forward
è 40 nell'esperimento 03 e 32 nel 02. L'encoder è identico fra le varianti dello
stesso esperimento.

L'encoder dell'esperimento~03 contiene \num{2660} parametri addestrabili,
identici nelle tre varianti e verificabili nel campo \code{transformer} di
\file{params.json}: per ciascuno dei due strati, \num{440} di auto-attenzione
($3\dvit^2 + 3\dvit$ per le proiezioni fuse più $\dvit^2 + \dvit$ per la
proiezione di uscita), \num{850} della rete feed-forward e $40$ delle due
\code{LayerNorm}, per un totale di \num{1330} per strato.

La rappresentazione del token \code{[CLS]} in uscita dall'encoder attraversa una
\code{LayerNorm} e una proiezione lineare sul numero di classi, seguita da
cross-entropy. Per PathMNIST, con nove classi, la testa ha
$10 \times 9 + 9 = 99$ parametri. La variante \code{vanilla} ne ha \num{4439} in
totale: \num{1480} nel \emph{patch embedding}, \num{2660} nell'encoder, 99 nella
testa e 200 fra token \code{[CLS]}, embedding posizionale e \code{LayerNorm}
finale. La scala ridotta contiene il costo di simulazione del circuito
(\cref{sec:overhead}) e circoscrive le conclusioni (\cref{sec:minacce}).

\section{Le due implementazioni}
\label{sec:implementazioni-vit}

Gli esperimenti 02 e 03 usano implementazioni distinte dello stesso schema
generale.

\begin{table}[htbp]
\centering
\caption[Confronto fra le due implementazioni di Vision Transformer]{%
Configurazioni delle due implementazioni di Vision Transformer.}
\label{tab:confronto-vit}
\small
\begin{tabularx}{\textwidth}{@{}l X X@{}}
\toprule
 & \classe{HybridQCNNViT} (esp.~02) & \classe{AblationViT} (esp.~03) \\
\midrule
Varianti & 2 (lineare / quantistica) & 3 (\code{vanilla}, \code{bounded\_mlp}, \code{quantum\_reg}) \\
\code{embed\_dim} & 8 & 10 \\
\code{ffn\_dim} & 32 & 40 \\
Init.\ \code{[CLS]} e posizionale & \code{randn} & \code{randn}$\,\times\,0{,}02$ \\
Conteggio per componente & no (solo totale) & sì (\code{count\_params}) \\
Metrica primaria & test accuracy & gap di generalizzazione \\
Parametri totali (PathMNIST) & \num{3169} / \num{3217} & \num{4439} / \num{4839} / \num{4499} \\
\bottomrule
\end{tabularx}
\end{table}

\subsection{Ambito dell'esperimento 02}
\label{ssec:precisazione-02}

Le architetture indicate in letteratura come Quantum Vision Transformer non
adottano un'unica definizione: alcune quantizzano blocchi di attenzione o MLP,
altre propongono trasformazioni quantistiche più estese
\parencite{cherrat2024quantumvision,comajoan2024qvt}. L'esperimento 02 considera
invece soltanto il \emph{patch embedding}.

Anche «quantum convolutional» può designare costruzioni differenti. La QCNN di
\textcite{cong2019qcnn} è un circuito gerarchico per dati quantistici, con
operazioni di convoluzione e pooling; la quanvolution di
\textcite{henderson2020quanvolutional} applica piccoli circuiti locali come
trasformazioni di patch classiche. Il modulo dell'esperimento 02 non coincide con
nessuna delle due famiglie: comprime ciascuna patch con uno strato lineare e
applica poi un VQC globale sugli otto valori ottenuti. La Self-Attention non è
quantistica, e l'unico componente quantizzato è il \emph{patch embedding}.

Gli esperimenti~02 e~03 usano la stessa struttura circuitale: compressione
lineare, riscalamento con $\tanh\cdot\pi$, \code{AngleEmbedding},
\code{StronglyEntanglingLayers} e misura di $\expval{Z_i}$ su ogni qubit. Nel 02
la compressione è interna a \classe{QuantumPatchEmbedding}; nel 03 è un
componente esplicito delle tre varianti. In entrambi i casi i suoi pesi sono
appaiati fra i modelli confrontati. Le due implementazioni non sono state
sottoposte a una prova di equivalenza numerica.
